\documentclass[spanish,letterpaper]{article}
\usepackage[utf8]{inputenc}
\usepackage[T1]{fontenc}
\usepackage[spanish,es-nodecimaldot]{babel}
\usepackage{helvet}
\renewcommand{\familydefault}{\sfdefault}
\usepackage[margin=1.65cm,top=1.2cm,bottom=1.45cm]{geometry}
\usepackage{amsmath,amssymb,array,tabularx,colortbl,graphicx,tikz,fancyhdr}
\usetikzlibrary{babel}
\usepackage{xcolor}

\definecolor{iiiepeNavy}{HTML}{163B5C}
\definecolor{iiiepeBlue}{HTML}{1F6F9F}
\definecolor{iiiepeAqua}{HTML}{20B7A8}
\definecolor{iiiepeOrange}{HTML}{F28C28}
\definecolor{iiiepePale}{HTML}{EEF5F8}
\definecolor{iiiepeInk}{HTML}{1E2A32}

\setlength{\parindent}{0pt}
\setlength{\parskip}{4pt}
\renewcommand{\arraystretch}{1.16}
\arrayrulecolor{iiiepeNavy!65}

\pagestyle{fancy}
\fancyhf{}
\renewcommand{\headrulewidth}{0pt}
\fancyfoot[L]{\scriptsize IIIEPE · MAEM · PFM02 · SD26}
\fancyfoot[C]{\scriptsize Portafolio de Evidencias · Bloque I}
\fancyfoot[R]{\scriptsize Pág. \thepage}

\begin{document}
\color{iiiepeInk}

\begin{center}
  \IfFileExists{img/iiiepe/marca-iiiepe.png}{\includegraphics[height=1.7cm]{img/iiiepe/marca-iiiepe.png}}{\textbf{\Large IIIEPE}}
  \vspace{0.25cm}

  {\large\bfseries Instituto de Investigación, Innovación y Estudios de Posgrado para la Educación del Estado de Nuevo León}\\[3pt]
  {\normalsize Maestría en Aprendizaje y Enseñanza de las Matemáticas (MAEM)}
  \vspace{0.35cm}

  \colorbox{iiiepeNavy}{%
    \parbox{0.92\textwidth}{%
      \centering\vspace{9pt}
      {\LARGE\bfseries\color{white} Portafolio de Evidencias · Bloque I}\\[4pt]
      {\normalsize\color{iiiepePale} Unidad de Aprendizaje: Temas Selectos de Matemáticas I (PFM02 · Sesiones 1--5)}
      \vspace{9pt}
    }%
  }
\end{center}

\vspace{0.25cm}
\noindent\begin{tabularx}{\textwidth}{|>{\bfseries\raggedright\arraybackslash}p{0.28\textwidth}|>{\raggedright\arraybackslash}X|}
\hline
\rowcolor{iiiepePale} Estudiante & Mtro. Martín Jonathan de la Cruz Muñoz \\\hline
Matrícula & 00874 \\\hline
\rowcolor{iiiepePale} Correo institucional & martin.delacruz@campus.iiiepe.edu.mx \\\hline
Programa & Maestría en Aprendizaje y Enseñanza de las Matemáticas \\\hline
\rowcolor{iiiepePale} Docente & Denis Concepción Escobedo Martínez \\\hline
Periodo / Cohorte & Ago--Nov 2026 (SD26) \\\hline
\end{tabularx}

\vspace{0.35cm}
{\large\bfseries\color{iiiepeNavy} Índice de Evidencias Integradas (Bloque I · Sesiones 1--5)}
\vspace{2pt}

Este volumen compila de manera ordenada los cinco reportes completos del Bloque I de \textit{Temas Selectos de Matemáticas I}, abarcando desde la conversión entre representaciones fraccionarias y decimales hasta la resolución de ecuaciones lineales y las propiedades de la igualdad:

\vspace{3pt}
\noindent\begin{tabularx}{\textwidth}{|>{\centering\arraybackslash}p{0.09\textwidth}|>{\raggedright\arraybackslash}X|>{\centering\arraybackslash}p{0.15\textwidth}|>{\centering\arraybackslash}p{0.13\textwidth}|}
\hline
\rowcolor{iiiepeNavy}
\color{white}\textbf{Act.} &
\color{white}\textbf{Título de la evidencia integrada} &
\color{white}\textbf{Extensión} &
\color{white}\textbf{Páginas} \\\hline
01 & \textbf{Conversión de fracciones a decimales y viceversa:} densidad, algoritmos de conversión y secuencias didácticas. & 9 pp. & 2--10 \\\hline
\rowcolor{iiiepePale}
02 & \textbf{Suma y resta con fracciones:} equivalencia, denominador común, recta numérica y retos contextualizados. & 7 pp. & 11--17 \\\hline
03 & \textbf{Multiplicación, división y números con signo:} modelos geométricos, inversión multiplicativa y reglas de signo. & 11 pp. & 18--28 \\\hline
\rowcolor{iiiepePale}
04 & \textbf{Proporcionalidad directa e inversa en un lanzador:} análisis de tablas, constante $k$ y variación funcional. & 6 pp. & 29--34 \\\hline
05 & \textbf{Ecuaciones lineales y propiedades de la igualdad:} balanza, propiedad uniforme y modelación algebraica. & 6 pp. & 35--40 \\\hline
\end{tabularx}

\vspace{0.3cm}
\noindent\colorbox{iiiepePale}{%
  \parbox{\dimexpr\textwidth-2\fboxsep\relax}{%
    \small\textbf{Nota de integración:} Cada actividad incluida en este portafolio conserva íntegramente su desarrollo matemático, representaciones gráficas, análisis de errores frecuentes y propuesta de intervención para el aula de matemáticas.
  }%
}

\end{document}
